% Generated by tools/textbook/export_listings.js from decks/plan.js and code/l05.S: do not edit.
\section{The program of this chapter}\label{sec:lab:l05}

The listings of this chapter are parts of one program, \fname{code/l05.S} in the
repository of the course, printed in full in Listing~\ref{lst:lab:l05}. The
automatic tests of the course material build it and run it on the headless
runner of \prgqt{}, with scripted input, and check the frame;
Figure~\ref{fig:lab:l05} is the frame that the test wrote.

\begin{description}[style=nextline,leftmargin=1.5em]
\item[On the screen.] The 32 bits of marker\_x as lamps. It starts at 2 (one lamp); each frame with LEFT held subtracts 1: after 2, 1 and 0 every lamp lights up, 0xFFFFFFFF, which is -1.
\item[What to press.] Hold LEFT for a few frames.
\end{description}

\begin{figure}[htb]
\centering
\includegraphics[width=0.5\textwidth]{lab_l05.png}
\caption{\fname{l05.S} on \prgqt{}.}\label{fig:lab:l05}
\end{figure}

\asmfile[caption={The complete program of the chapter, \fname{code/l05.S}.},label=lst:lab:l05]{code/l05/program.S}

To run it on \prgqt{} the commands are those of Listing~\ref{lst:lab:l05:run};
Section~\ref{sec:prg:testing} gives the same steps for QEMU and for the board.

\begin{shell}[caption={Building \fname{l05.S} and running it on \prgqt{}.},label=lst:lab:l05:run]
# 1. start PRG32-QT on this computer (it listens on port 8080)

# 2. in the PRG32 folder, with ESP-IDF active: build the cartridge
python3 -m prg32 cartridge build ~/arc12/code/l05.S --portable \
    --entry-prefix l05 --name l05 --out l05.prg32

# 3. send it to PRG32-QT and start it
python3 -m prg32 esp32c6 upload-and-run l05.prg32 \
    --url http://localhost:8080 --slot cart0

# 4. look, press, compare with the previous slide
\end{shell}
